\documentclass[border=12pt]{standalone}
\usepackage[utf8]{inputenc}
\usepackage[T1]{fontenc}
\usepackage[spanish]{babel}
\usepackage{tikz}
\usetikzlibrary{positioning, arrows.meta, backgrounds, shapes.geometric}

\begin{document}
\begin{tikzpicture}[
    font=\small,
    nano/.style   = {draw, thick, rounded corners, fill=blue!12,
                     minimum width=3.2cm, minimum height=1.2cm, align=center},
    esp/.style    = {draw, thick, rounded corners, fill=violet!14,
                     minimum width=3.2cm, minimum height=1.1cm, align=center},
    mpr/.style    = {draw, thick, rounded corners, fill=green!14,
                     minimum width=3.2cm, minimum height=0.95cm, align=center},
    perif/.style  = {draw, fill=gray!8, align=left, text width=4.8cm,
                     inner sep=5pt, rounded corners},
    cloud/.style  = {draw, thick, ellipse, fill=orange!18,
                     minimum width=3.2cm, minimum height=1.4cm, align=center},
    bus/.style    = {line width=1.4pt},
    wire/.style   = {thick},
    uart/.style   = {thick, {Latex[length=2mm]}-{Latex[length=2mm]}}
]

% ===================== BUS I2C =====================
\draw[bus] (2,6.5) -- (17.5,6.5);
\node[above=1pt] at (9.75,6.55)
    {\textbf{BUS I\textsuperscript{2}C}\quad SDA = A4\;\textbullet\;SCL = A5\;\textbullet\;pull-ups 4.7\,k$\Omega$ a +5V};
% pull-up simbolico a +5V
\draw[wire] (2.7,6.5) -- (2.7,7.5) node[above, font=\footnotesize] {+5V};
\node[font=\footnotesize] at (3.35,7.0) {4.7k};

% ===================== NODOS SOBRE EL BUS =====================
\node[nano] (master) at (4.5,4.5)  {\textbf{MAESTRO}\\ Arduino Nano\\ (FSM del proceso)};
\node[nano] (slave1) at (10.5,4.5) {\textbf{ESCLAVO 1}\\ Arduino Nano\\ 0x30};
\node[nano] (slave2) at (16.5,4.5) {\textbf{ESCLAVO 2}\\ Arduino Nano\\ 0x31};

% MPR121 colgado del bus
\node[mpr] (mpr) at (7.5,8) {MPR121 táctil\\ 0x5A (level-shifter)};
\draw[wire] (7.5,6.5) -- (mpr.south);

% conexiones nodo -> bus (SDA/SCL)
\foreach \x in {4.5, 10.5, 16.5}{
    \draw[wire] (\x,5.1) -- (\x,6.5) node[midway,right,font=\footnotesize] {SDA/SCL};
}

% ===================== ESP32 + NUBE (UART / WiFi) =====================
\node[esp]   (esp)  at (-0.5,2.5) {\textbf{ESP32} DevKit v1\\ (WiFi + Adafruit)};
\node[cloud] (nube) at (-0.5,5.3) {\textbf{Adafruit IO}\\ Dashboard};
\draw[uart] (esp.north) -- (master.west) node[midway,above,sloped,font=\footnotesize] {UART D0/D1};
\draw[thick, dashed, {Latex}-{Latex}] (esp.north) -- (nube.south)
      node[midway, right, font=\footnotesize] {WiFi};

% ===================== PERIFERICOS (paneles bajo cada nodo) =====================
\node[perif] (mp) at (4.5,1.35) {%
    \textbf{Periféricos del maestro}\\[2pt]
    LCD 2x16 (8 bits): D2--D12\\
    Servo 360 giro continuo: D13\\
    Botón START: A2\\
    LED táctil: A1};

\node[perif] (s1p) at (10.5,1.05) {%
    \textbf{Esclavo 1 -- corredera}\\[2pt]
    IR presencia: D2\\
    Ultrasónico: TRIG D3 / ECHO D4\\
    Motor DC (L298): IN1 D5 / IN2 D6\\
    Servo corredera: D11\\
    Servo tapadera: D10};

\node[perif] (s2p) at (16.5,1.05) {%
    \textbf{Esclavo 2 -- dispensado}\\[2pt]
    HW-503 temperatura: A1\\
    Stepper condimento (ULN2003): D2--D5\\
    Stepper chile (ULN2003): D6--D9\\
    Servo cafetera (giro continuo): D10\\
    Relé: D11};

% lineas nodo -> panel
\draw[wire] (master.south) -- (mp.north);
\draw[wire] (slave1.south) -- (s1p.north);
\draw[wire] (slave2.south) -- (s2p.north);

% ===================== NOTA =====================
\node[align=left, font=\footnotesize, anchor=north west] at (-1.5,-1.6) {%
    Notas: GND común a todos los nodos. Servomotores y motor DC con fuente de 5V\\
    independiente (GND común). El maestro concentra la lógica; los esclavos solo\\
    reportan sensores y ejecutan comandos vía el mapa de registros I\textsuperscript{2}C.};

\end{tikzpicture}
\end{document}
